% 09_run.tex - run the coupled model and read the summary
\section{Run the coupled model}

\step{Run Raven}
\term{t2_run.txt}
\see \msg{Exiting Gracefully: Successful Simulation}. The run takes about \V{t2:runsec}~seconds. The line
\cmd{Groundwater steady-state initialization converged in \V{t2:ssiter} outer iterations} confirms the steady start.
\begin{tocheck}
Raven prints \msg{Exiting Gracefully: Successful Simulation}. The command
\cmd{grep ERROR out/Raven\_errors.txt} prints nothing.
\end{tocheck}
\begin{trouble}
\begin{itemize}
\item \emph{``HRU \dots has an AQUIFER\_PROFILE but no polygon''}: the HRU numbers in \cmd{hrus.geojson} and in
  \cmd{Creek.rvh} do not match, or the field is not called \cmd{HRU\_ID}.
\item \emph{``\dots does not overlap any active groundwater cell''}: a small HRU under large cells. Make
  \cmd{:GridCellSize} or \cmd{:MinCellCoverage} smaller.
\item \emph{``:WellFlowRate 1 in the .rvt file names no well''} or \emph{``:BoundaryHead POND \dots names no
  :GeneralHeadBoundary or :SpecifiedHeadBoundary''}: the ID or the name differs between the \cmd{.rvt} and the
  \cmd{.rvg} file.
\item A missing \cmd{:InputCRS} line: see Step~\ref{st:nocrs}.
\item Compare your files with the finished ones: \cmd{diff -wB \$WS/model/finished/Creek.rvg Creek.rvg} prints the lines that
  differ, leaving spacing and empty lines aside (nothing if the files are the same).
\end{itemize}
\end{trouble}

\step{Read what Raven built}
\term{t2_summary.txt}
\begin{tocheck}
\begin{itemize}
\item The grid has \V{t2:nlay}~layers of \V{t2:nrow}$\times$\V{t2:ncol} cells and \V{t2:activecols} active columns. The
  \V{t2:inactivecols} other columns lie under the ridge or outside the HRUs. \V{t2:fidelity}\% of the active area lies in
  the HRU that covers most of each cell.
\item There are \V{t2:rivcells} river cells, one well cell, \V{t2:ghbcells} GHB cells (in both layers) and
  \V{t2:chdcells} CHD cells under the pond.
\item The observation wells are in layer \V{t2:owlay101} (OW\_NEAR), layer \V{t2:owlay102} (OW\_UPPER) and layer
  \V{t2:owlay103} (OW\_SLOPE; the top layer is dry there).
\item In the HRU table, \cmd{linked\_fraction\_of\_polygon} is close to 1 for every HRU. A low value means that an HRU lies
  outside the grid, or that a projection is wrong.
\end{itemize}
\end{tocheck}
Figure~\ref{fig:t2grid} shows the active cells (from \cmd{out/mf6/grid\_cells.geojson}) over the HRUs.
\fig[0.86]{t2_grid.png}{The active cells of the grid that Raven built, over the HRUs. There are none under the ridge
(HRU~4), which has no aquifer.\label{fig:t2grid}}
